\documentclass[aspectratio=169]{beamer}

% ============================================
% Theme and Colors
% ============================================
\usetheme{moloch}
\molochset{
    block=fill,
    background=dark,
    sectionpage=none
}
\usecolortheme{dracula}

% ============================================
% Color Mapping & Dracula Overrides
% ============================================
\definecolor{birdgreen}{RGB}{80, 250, 123}  % draculagreen
\definecolor{skyblue}{RGB}{139, 233, 253}    % draculacyan
\definecolor{darkblue}{RGB}{40, 42, 54}      % draculabg

\setbeamercolor{title separator}{fg=draculapink}
\setbeamercolor{progress bar}{fg=draculagreen, bg=draculacurrent}

% ============================================
% TikZ Styles (Updated for Dracula Dark Mode)
% ============================================
\tikzset{
    sedbox/.style={rectangle, draw=draculapurple, fill=draculacurrent, rounded corners=5pt, line width=1.2pt, minimum width=6.5cm, minimum height=1cm, align=center, font=\small, text=draculafg},
    greenbox/.style={rectangle, draw=draculagreen, fill=draculacurrent, rounded corners=5pt, line width=1.2pt, minimum width=6.5cm, minimum height=1cm, align=center, font=\small, text=draculafg},
    orangebox/.style={rectangle, draw=draculaorange, fill=draculacurrent, rounded corners=5pt, line width=1.2pt, minimum width=6.5cm, minimum height=1cm, align=center, font=\small, text=draculafg},
    sednote/.style={rectangle, draw=draculacomment, fill=draculabg, rounded corners=5pt, line width=0.8pt, minimum width=3cm, align=left, font=\tiny, text=draculafg},
    sedarrow/.style={-Stealth, thick, color=draculacyan, line width=1.2pt}
}

% ============================================
% Packages
% ============================================
\usepackage[utf8]{inputenc}
\usepackage[T1]{fontenc}
\usepackage{graphicx}
\graphicspath{{images/}}
\usepackage{tikz}
\usetikzlibrary{shapes.geometric, arrows.meta, positioning, calc, decorations.pathreplacing}
\usepackage{booktabs}
\usepackage{colortbl}
\usepackage{amsmath}
\usepackage{xcolor}
\usepackage{hyperref}
\usepackage{listings}
\usepackage{multicol}


% ============================================
% Listing Style
% ============================================
\lstset{
    basicstyle=\tiny\ttfamily,
    keywordstyle=\color{blue},
    commentstyle=\color{gray},
    breaklines=true,
    frame=single
}

% ============================================
% Title Information
% ============================================
\title[Species Identification in Noisy Soundscapes]{\textbf{Species Identification in Noisy Soundscapes}}
\subtitle{A Deep Learning Approach using Sound Event Detection}
\author{Gourove Roy(2105017), Niloy Das Robin(2105019),\\ Sayjad Rahman(2105021)}
\institute{CSE 329 - Machine Learning Sessional\\BUET}
\date{\textcolor{draculapink}{April 2026}}

% ============================================
% Document Begin
% ============================================
\begin{document}

% ------------------------------------------
% Title Slide
% ------------------------------------------
\begin{frame}
    \titlepage
\end{frame}

% ------------------------------------------
% Outline
% ------------------------------------------
\begin{frame}{Outline}
    \tableofcontents
\end{frame}

% ============================================
\section{Introduction}
% ============================================

% ------------------------------------------
% Problem Statement
% ------------------------------------------
\begin{frame}{Problem Statement}
    \begin{columns}
        \begin{column}{0.55\textwidth}
            \begin{block}{Acoustic Species Identification}
                Identify bird, insect, and amphibian species from \textbf{noisy soundscape recordings} in tropical rainforest environments.
            \end{block}

            \vspace{0.3cm}
            \textbf{Key Challenges:}
            \begin{itemize}
                \item 206 different species to classify
                \item Heavy background noise (rain, wind, other animals)
                \item Overlapping species calls
                \item Extreme class imbalance (some species very rare)
                \item 60-second long soundscape recordings
            \end{itemize}
        \end{column}
        \begin{column}{0.42\textwidth}
            \begin{tikzpicture}[scale=0.8]
                % Soundscape representation
                \draw[draw=draculapurple, fill=draculacurrent, rounded corners, line width=1pt] (0,0) rectangle (5,3.5);
                \node[text=draculafg] at (2.5,3.8) {\textbf{60s Soundscape}};

                % Multiple species
                \draw[draw=draculapurple, fill=draculagreen!60, rounded corners] (0.2,2.5) rectangle (2,3);
                \node[font=\tiny, text=draculabg] at (1.1,2.75) {Bird A};

                \draw[draw=draculapurple, fill=draculaorange!60, rounded corners] (1.5,1.8) rectangle (3.5,2.3);
                \node[font=\tiny, text=draculabg] at (2.5,2.05) {Insect B};

                \draw[draw=draculapurple, fill=draculapink!60, rounded corners] (2.8,1) rectangle (4.8,1.5);
                \node[font=\tiny, text=draculabg] at (3.8,1.25) {Frog C};

                % Noise overlay
                \draw[draw=draculacomment, fill=draculacomment, opacity=0.5] (0,0) rectangle (5,0.7);
                \node[font=\tiny, text=draculafg] at (2.5,0.35) {Background Noise};
            \end{tikzpicture}
        \end{column}
    \end{columns}
\end{frame}

% ------------------------------------------
% Dataset Overview
% ------------------------------------------
\begin{frame}{Dataset Overview}
    \begin{columns}[T, onlytextwidth]
        \begin{column}{0.46\textwidth}
            \begin{block}{Training Data}
                \small
                \begin{itemize}
                    \item \textbf{Species}: 206 total
                    \begin{itemize}
                        \item Aves (birds): majority
                        \item Amphibia: 6 species
                        \item Insecta: multiple species
                    \end{itemize}
                    \item \textbf{Audio format}: .ogg files
                    \item \textbf{Sample rate}: 32,000 Hz
                    \item \textbf{Duration}: Variable (5s - 60s)
                \end{itemize}
            \end{block}

            \begin{alertblock}{Class Imbalance Problem}
                \small Some species have 1000+ samples while others have $<$ 10 samples
            \end{alertblock}
        \end{column}
        \begin{column}{0.46\textwidth}
            \begin{block}{Test Data (Soundscapes)}
                \small
                \begin{itemize}
                    \item 60-second continuous recordings
                    \item Multiple species per recording
                    \item Real-world noise conditions
                    \item Predictions required per 5-second segment
                \end{itemize}
            \end{block}

            \vspace{0.1cm}
            \textbf{Evaluation Metric:}
            \small
            \begin{itemize}
                \item Row-wise micro-averaged F1
                \item Top-1 and Top-5 accuracy
            \end{itemize}
        \end{column}
    \end{columns}
\end{frame}

% ============================================
\section{Audio Processing Pipeline}
% ============================================

\begin{frame}{Audio Processing Pipeline}
    \begin{center}
        \resizebox{\textwidth}{!}{
            \begin{tikzpicture}[
                node distance=0.8cm,
                arrow/.style={sedarrow}
            ]
                % Nodes
                \node[sedbox, minimum width=2.2cm, minimum height=0.8cm] (audio) {\textbf{Raw Audio}\\(20s, 32kHz)};
                \node[sedbox, right=of audio, minimum width=2.2cm, minimum height=0.8cm] (mel) {\textbf{Mel Spectrogram}\\(224 $\times$ 512)};
                \node[sedbox, right=of mel, minimum width=2.2cm, minimum height=0.8cm] (norm) {\textbf{Normalization}\\(Z-score + MinMax)};
                \node[sedbox, right=of norm, minimum width=2.2cm, minimum height=0.8cm] (chan) {\textbf{3-Channel}\\Replication};
                \node[sedbox, right=of chan, minimum width=2.2cm, minimum height=0.8cm] (cnn) {\textbf{CNN Input}\\(3 $\times$ 224 $\times$ 512)};

                % Arrows
                \draw[arrow] (audio) -- (mel);
                \draw[arrow] (mel) -- (norm);
                \draw[arrow] (norm) -- (chan);
                \draw[arrow] (chan) -- (cnn);
            \end{tikzpicture}
        }
    \end{center}

    \vspace{0.2cm}
    \begin{columns}[T, onlytextwidth]
        \begin{column}{0.46\textwidth}
            \begin{block}{Mel Spectrogram Parameters}
                \small
                \begin{itemize}
                    \item \textbf{Sample Rate}: 32,000 Hz
                    \item \textbf{Mel Bands}: 224
                    \item \textbf{n\_fft}: 4,096
                    \item \textbf{Hop Length}: 1,252
                    \item \textbf{Freq Range}: 0 - 16 kHz
                \end{itemize}
            \end{block}
        \end{column}
        \begin{column}{0.46\textwidth}
            \begin{block}{Why 20-Second Chunks?}
                \small
                \begin{itemize}
                    \item Captures full Amphibia calls
                    \item Provides temporal context
                    \item Optimal accuracy vs. time
                    \item 20s (87.2\%) vs 5s (84.2\%)
                \end{itemize}
            \end{block}
        \end{column}
    \end{columns}
\end{frame}

% ============================================
\section{Model Architecture}
% ============================================

\begin{frame}{Sound Event Detection (SED) Architecture}
    \begin{center}
        \resizebox{!}{0.75\textheight}{
            \begin{tikzpicture}[node distance=0.4cm]
                % Main nodes
                \node[sedbox] (input) {\textbf{1. Audio Input}\\20-second waveform sampled at 32 kHz};
                \node[sedbox, below=of input] (spectro) {\textbf{2. Spectrogram Front-End}\\Convert waveform into Mel spectrogram\\224 mel bands, 4096 n\_fft};
                \node[sedbox, below=of spectro] (norm) {\textbf{3. Normalization}\\Absmax normalization on audio\\then z-score + min-max on the spectrogram};
                \node[sedbox, below=of norm] (rep) {\textbf{4. 3-Channel Replication}\\Turn the spectrogram into an image-like tensor};
                \node[sedbox, below=of rep] (backbone) {\textbf{5. CNN Backbone}\\Use a features\_only backbone to preserve time structure\\Examples: EfficientNet, NFNet, RegNet};
                \node[sedbox, below=of backbone] (head) {\textbf{6. SED Head}\\Apply GeMFreq across frequency, then AttHead\\Dense + ReLU + Dropout + Conv1d for framewise logits};

                % Arrows
                \draw[sedarrow] (input) -- (spectro);
                \draw[sedarrow] (spectro) -- (norm);
                \draw[sedarrow] (norm) -- (rep);
                \draw[sedarrow] (rep) -- (backbone);
                \draw[sedarrow] (backbone) -- (head);

                % Side notes
                \node[sednote, left=1cm of head] (backbones) {\textbf{Backbone Examples}\\\textbullet\ tf\_efficientnet\_b3 / b0\\\textbullet\ eca\_nfnet\_l0\\\textbullet\ regnety\_016 / 008};
                \node[sednote, right=1cm of head] (outputs) {Framewise output: [B, 206, T]\\Clip output: max over time $\rightarrow$ [B, 206]};
                
                \draw[sedarrow, dashed, thin] (backbones) -- (head);
                \draw[sedarrow, dashed, thin] (head) -- (outputs);
            \end{tikzpicture}
        }
    \end{center}
\end{frame}

% ------------------------------------------
% GeMFreq and AttHead Details
% ------------------------------------------
\begin{frame}[t]{Key Architecture Components}
    \begin{columns}[T, onlytextwidth]
        \begin{column}{0.46\textwidth}
            \begin{block}{GeMFreq Pooling}
                \small
                \textbf{Generalized Mean Freq Pooling}
                \begin{itemize}
                    \item Learnable pooling parameter $p$
                    \item Formula: $\left(\frac{1}{|F|}\sum_{f} x_f^p\right)^{1/p}$
                    \item Adapts: max vs average pooling
                    \item Captures frequency patterns
                \end{itemize}
            \end{block}

            \begin{exampleblock}{Why Features Only?}
                \small Spatial features w/o classification head $\rightarrow$ preserves temporal dimension
            \end{exampleblock}
        \end{column}
        \begin{column}{0.46\textwidth}
            \begin{block}{AttHead (Attention Head)}
                \small
                \begin{enumerate}
                    \item Dense: backbone\_dim $\rightarrow$ 512
                    \item Parallel branches:
                    \begin{itemize}
                        \item \textbf{Attention}: 1$\times$1 Conv $\rightarrow$ Softmax
                        \item \textbf{Logits}: 1$\times$1 Conv $\rightarrow$ 206
                    \end{itemize}
                    \item Weighted aggregation:
                    $$\text{clip\_logits} = \sum_t \text{att}_t \cdot \text{logits}_t$$
                \end{enumerate}
            \end{block}

            \small \textbf{Advantage}: Learns which time frames contain species calls!
        \end{column}
    \end{columns}
\end{frame}

% ============================================
\section{Training Techniques}
% ============================================

% ------------------------------------------
% Data Augmentation
% ------------------------------------------
\begin{frame}[t]{Data Augmentation Techniques}
    \begin{columns}[T]
        \begin{column}{0.32\textwidth}
            \begin{block}{Mixup ($\alpha = 0.5$)}
                \begin{center}
                    \begin{tikzpicture}[scale=0.55]
                        \draw[draw=draculacyan, fill=draculacyan!60] (0,0) rectangle (2,1.5);
                        \node[text=draculabg] at (1,0.75) {\tiny Audio A};
                        \draw[draw=draculapink, fill=draculapink!60] (0,1.8) rectangle (2,3.3);
                        \node[text=draculabg] at (1,2.55) {\tiny Audio B};
                        \draw[sedarrow, color=draculafg] (2.2,1.65) -- (3,1.65);
                        \draw[draw=draculapurple, fill=draculapurple!60] (3.2,0.8) rectangle (5.2,2.5);
                        \node[text=draculabg] at (4.2,1.65) {\tiny Mixed};
                    \end{tikzpicture}
                \end{center}
                \small $\tilde{x} = \lambda x_a + (1-\lambda) x_b$

                \textcolor{draculagreen}{\textbf{+2-3\% accuracy}}
            \end{block}
        \end{column}
        \begin{column}{0.32\textwidth}
            \begin{block}{SpecAugment}
                \begin{center}
                    \begin{tikzpicture}[scale=0.55]
                        \draw[draw=draculapurple, fill=draculapurple!40] (0,0) rectangle (3,2);
                        \draw[fill=draculagreen] (1,0) rectangle (1.3,2);
                        \draw[fill=draculagreen] (0,0.8) rectangle (3,1.1);
                        \node[text=draculagreen] at (1.5,-0.3) {\tiny Masked};
                    \end{tikzpicture}
                \end{center}
                \begin{itemize}
                    \item \small Time: 20\%
                    \item \small Freq: 10\%
                \end{itemize}
                \textcolor{draculagreen}{\textbf{+1-2\% accuracy}}
            \end{block}
        \end{column}
        \begin{column}{0.32\textwidth}
            \begin{block}{Weighted Sampling}
                \small Address imbalance:
                $$w_c = \frac{1}{\sqrt{n_c}}$$
                \begin{itemize}
                    \item \small Rare species $\uparrow$
                    \item \small Common $\downarrow$
                \end{itemize}
                \textcolor{draculagreen}{\textbf{Critical for rare}}
            \end{block}
        \end{column}
    \end{columns}
\end{frame}

% ------------------------------------------
% Training Configuration
% ------------------------------------------
\begin{frame}[t]{Training Configuration}
    \begin{columns}[T, onlytextwidth]
        \begin{column}{0.46\textwidth}
            \begin{block}{Optimizer \& Scheduler}
                \footnotesize
                \begin{itemize}
                    \item \textbf{Optimizer}: AdamW
                    \begin{itemize}
                        \item LR: $5 \times 10^{-4}$
                        \item WD: $10^{-4}$
                    \end{itemize}
                    \item \textbf{Scheduler}: CosineAnnealing
                    \begin{itemize}
                        \item $T_0 = 5$ epochs
                        \item $\eta_{min} = 10^{-6}$
                    \end{itemize}
                    \item \textbf{Grad Clip}: max\_norm = 1.0
                \end{itemize}
            \end{block}
        \end{column}
        \begin{column}{0.46\textwidth}
            \begin{block}{Loss Function}
                \footnotesize
                \textbf{CrossEntropy + Label Smoothing}
                \vspace{-0.1cm}
                $$\mathcal{L} = -\sum_c y_c' \log(p_c)$$
                \vspace{-0.3cm}
                \begin{itemize}
                    \item Smoothing: $\epsilon = 0.1$
                    \item Prevents overconfidence
                    \item Better calibration
                \end{itemize}
            \end{block}

            \begin{exampleblock}{Training Details}
                \footnotesize
                \begin{itemize}
                    \item Epochs: 15-35
                    \item Batch size: 64
                    \item 5-fold cross-validation
                \end{itemize}
            \end{exampleblock}
        \end{column}
    \end{columns}
\end{frame}

\begin{frame}{Supervised Training Pipeline}
    \begin{center}
        \resizebox{!}{0.75\textheight}{
            \begin{tikzpicture}[node distance=0.4cm]
                % Main nodes
                \node[greenbox] (ds) {\textbf{1. Labeled Dataset}\\Prepare 20-second chunks with species labels};
                \node[greenbox, below=of ds] (aug) {\textbf{2. Training Augmentation}\\Apply MixUp and SpecAugment during training only\\This improves robustness to noisy audio conditions};
                \node[greenbox, below=of aug] (samp) {\textbf{3. Weighted Sampling}\\Oversample rare species to reduce class imbalance};
                \node[greenbox, below=of samp] (pass) {\textbf{4. Forward Pass}\\Run the SED model and convert framewise evidence\\into clip-level logits for supervised learning};
                \node[greenbox, below=of pass] (loss) {\textbf{5. Loss and Optimization}\\CrossEntropy with label smoothing, AdamW optimizer,\\cosine learning-rate schedule, and gradient clipping};
                \node[greenbox, below=of loss] (best) {\textbf{6. Checkpoint Selection}\\Save the best validation model as the base checkpoint};

                % Arrows
                \draw[sedarrow] (ds) -- (aug);
                \draw[sedarrow] (aug) -- (samp);
                \draw[sedarrow] (samp) -- (pass);
                \draw[sedarrow] (pass) -- (loss);
                \draw[sedarrow] (loss) -- (best);

                % Side notes
                \node[sednote, right=1cm of loss] (settings) {\textbf{Key Settings}\\\textbullet\ MixUp alpha = 0.5\\\textbullet\ Label smoothing = 0.1\\\textbullet\ AdamW lr = 1e-4\\\textbullet\ Weight decay = 1e-2};
                \draw[sedarrow, dashed, thin] (loss) -- (settings);
            \end{tikzpicture}
        }
    \end{center}
\end{frame}

% ============================================
\section{Inference Pipeline}
% ============================================

\begin{frame}{Overlap-Average Inference}
    \centering
    \resizebox{!}{0.5\textheight}{
        \begin{tikzpicture}[scale=0.8]
            % 60s soundscape
            \draw[draw=draculapurple, fill=draculacurrent, rounded corners, line width=1.2pt] (0,3) rectangle (12,3.8);
            \node[text=draculafg] at (6,3.4) {\textbf{60-second Soundscape}};

            % Segments
            \foreach \i in {0,1,2,3} {
                \pgfmathsetmacro{\start}{\i*2}
                \pgfmathsetmacro{\midy}{2.2 - \i*0.5}
                \draw[draw=draculacyan, fill=draculagreen!\i 0!draculapurple, opacity=0.5, rounded corners, line width=0.8pt]
                    (\start,\midy) rectangle (\start+4,\midy+0.4);
                \node[font=\tiny, text=draculafg] at (\start+2,\midy+0.2) {Chunk \pgfmathparse{int(\i+1)}\pgfmathresult};
            }

            % Overlap indication
            \draw[<->, thick, draculared] (2,1.4) -- (4,1.4);
            \node[font=\tiny, draculared] at (3,1.1) {Overlap};

            % Arrow to averaging
            \draw[sedarrow] (6,0.5) -- (6,-0.1);

            % Averaged predictions
            \node[sedbox, fill=draculacurrent, draw=draculapink, minimum width=8cm, minimum height=0.8cm] at (6,-0.6) {\textbf{Overlap-Averaged Framewise Predictions}};

            % Final output
            \draw[sedarrow] (6,-1.1) -- (6,-1.5);
            \node[sedbox, fill=draculacurrent, draw=draculaorange, minimum width=6cm, minimum height=0.8cm] at (6,-2.0) {\textbf{12 $\times$ 206 Prediction Matrix}};
        \end{tikzpicture}
    }

    \begin{columns}[T, onlytextwidth]
        \begin{column}{0.46\textwidth}
            \footnotesize \textbf{Process:}
            \begin{enumerate}
                \item Split into 12 overlapping chunks
                \item Chunk $\rightarrow$ framewise predictions
                \item Average overlapping frame preds
            \end{enumerate}
        \end{column}
        \begin{column}{0.46\textwidth}
            \footnotesize \textbf{Benefits:}
            \begin{itemize}
                \item More robust predictions
                \item Implicit TTA effect
                \item +0.2-0.3\% accuracy boost
            \end{itemize}
        \end{column}
    \end{columns}
\end{frame}

\begin{frame}[t]{Delta-Shift TTA \& Gaussian Smoothing}
    \begin{columns}[T, onlytextwidth]
        \begin{column}{0.46\textwidth}
            \begin{block}{Delta-Shift TTA}
                \small
                Temporal shift augmentation:
                \begin{center}
                    \resizebox{\textwidth}{!}{
                        \begin{tikzpicture}[node distance=0.2cm]
                            \node[sedbox, fill=draculacurrent, draw=draculacyan, minimum width=4cm, minimum height=0.5cm] (orig) {\tiny Original};
                            \node[sedbox, fill=draculacurrent, draw=draculacyan!70, minimum width=4cm, minimum height=0.5cm, below=of orig, xshift=-0.3cm] (s1) {\tiny Shift -2};
                            \node[sedbox, fill=draculacurrent, draw=draculacyan!70, minimum width=4cm, minimum height=0.5cm, below=of s1, xshift=0.6cm] (s2) {\tiny Shift +2};

                            \node[sedbox, fill=draculagreen!20, draw=draculagreen, right=1.2cm of s1, minimum width=1.8cm, minimum height=0.8cm] (blend) {\textbf{\textcolor{draculabg}{Blend}}};

                            \draw[sedarrow] (orig.east) -- (blend.north west);
                            \draw[sedarrow] (s1.east) -- (blend.west);
                            \draw[sedarrow] (s2.east) -- (blend.south west);
                        \end{tikzpicture}
                    }
                \end{center}
                \vspace{-0.2cm}
                \textbf{Blending weights:}
                $$\text{pred} = 0.50 \cdot p_0 + 0.25 \cdot p_{\pm 2}$$
            \end{block}
        \end{column}
        \begin{column}{0.46\textwidth}
            \begin{block}{Gaussian Smoothing}
                \small
                1D convolution to reduce jitter:
                \textbf{Kernel:} $[0.1, 0.2, 0.4, 0.2, 0.1]$
                \begin{center}
                    \begin{tikzpicture}[scale=0.5]
                        % Before
                        \draw[thick, draculacyan] plot[smooth] coordinates {
                            (0,1) (0.5,2) (1,1.2) (1.5,2.5) (2,1.8) (2.5,2.2) (3,1.5)
                        };
                        \node[text=draculafg] at (1.5,-0.3) {\tiny Before};
                        % Arrow
                        \draw[sedarrow] (3.5,1.5) -- (4.5,1.5);
                        % After
                        \draw[thick, draculagreen] plot[smooth] coordinates {
                            (5,1.3) (5.5,1.7) (6,1.8) (6.5,2) (7,1.9) (7.5,1.8) (8,1.6)
                        };
                        \node[text=draculafg] at (6.5,-0.3) {\tiny After};
                    \end{tikzpicture}
                \end{center}
                Reduces noise while preserving patterns.
            \end{block}
        \end{column}
    \end{columns}

    \vspace{0.1cm}
    \begin{exampleblock}{Combined Effect}
        \small Delta-Shift TTA + Gaussian Smoothing = \textbf{Free accuracy boost!}
    \end{exampleblock}
\end{frame}

\begin{frame}{Inference \& Ensemble Pipeline}
    \begin{center}
        \resizebox{!}{0.75\textheight}{
            \begin{tikzpicture}[node distance=0.4cm]
                % Main nodes
                \node[orangebox] (input) {\textbf{1. Soundscape Input}\\Start from a 60-second audio recording};
                \node[orangebox, below=of input] (win) {\textbf{2. Windowing}\\Split the recording into overlapping 20-second chunks\\using a 5-second step size};
                \node[orangebox, below=of win] (pred) {\textbf{3. Per-Chunk SED Prediction}\\Run each chunk through the SED model to get framewise probabilities};
                \node[orangebox, below=of pred] (refine) {\textbf{4. Temporal Refinement}\\Combine overlap-average inference, Delta-Shift TTA,\\and Gaussian smoothing for more stable predictions};
                \node[orangebox, below=of refine] (scores) {\textbf{5. Segment-Level Scores}\\Max-pool within each 5-second segment};

                % Arrows
                \draw[sedarrow] (input) -- (win);
                \draw[sedarrow] (win) -- (pred);
                \draw[sedarrow] (pred) -- (refine);
                \draw[sedarrow] (refine) -- (scores);

                % Side notes and ensemble
                \node[sednote, right=1.2cm of pred] (ens) {\textbf{Ensemble Members}\\\textbullet\ tf\_efficientnet\_b3\\\textbullet\ tf\_efficientnet\_b0\\\textbullet\ eca\_nfnet\_l0\\\textbullet\ tf\_efficientnet\_b4\\\textbullet\ regnety\_016 / 008};
                
                \node[orangebox, right=1.2cm of scores] (mean) {\textbf{6. Arithmetic Mean}\\Average predictions across all models};
                \node[orangebox, below=0.8cm of mean] (out) {\textbf{7. Final Output}\\Final prediction per 5-second segment};

                % Connection arrows
                \draw[sedarrow] (scores.east) -- (mean.west);
                \draw[sedarrow] (mean) -- (out);
                \draw[sedarrow, dashed, thin] (ens) -- (mean);
            \end{tikzpicture}
        }
    \end{center}
\end{frame}

\begin{frame}{Enabling Multi-Label Classification}
    \begin{block}{From Single-Label Training to Multi-Label Inference}
        \footnotesize
        Although trained on \textbf{primary labels} only, the model effectively identifies multiple concurrent species through its architectural design:
    \end{block}

    \begin{columns}[T, onlytextwidth]
        \begin{column}{0.46\textwidth}
            \begin{exampleblock}{Architectural Support}
                \footnotesize
                \begin{itemize}
                    \item \textbf{Sigmoid Activation}: Used at inference instead of Softmax. Each class is treated as an independent binary detection task.
                    \item \textbf{SED Localization}: The model learns features for specific species calls in time-frequency bins, allowing evidence for multiple species to exist simultaneously.
                \end{itemize}
            \end{exampleblock}
        \end{column}
        \begin{column}{0.46\textwidth}
            \begin{exampleblock}{Inference Logic}
                \footnotesize
                \begin{itemize}
                    \item \textbf{Class Independence}: A high score for one species does not force low scores for others (unlike Softmax).
                    \item \textbf{Thresholding}: Any species score exceeding a pre-defined threshold (e.g., 0.3) is included in the final prediction list for that 5s segment.
                \end{itemize}
            \end{exampleblock}
        \end{column}
    \end{columns}
    
    \vspace{0.2cm}
    \centering \textcolor{draculapurple}{\textbf{Result: One model can detect 0, 1, or many species in a single audio chunk.}}
\end{frame}

% ============================================
\section{Multi-Model Ensemble}
% ============================================

% ------------------------------------------
% Ensemble Architecture
% ------------------------------------------
\begin{frame}{6-Model Ensemble Architecture}
    \begin{center}
        \begin{tabular}{llcc}
            \toprule
            \textbf{Model} & \textbf{Backbone} & \textbf{Feature Dim} & \textbf{Training Stage} \\
            \midrule
            1 & tf\_efficientnet\_b3 & 384 & Trained (ours) \\
            2 & tf\_efficientnet\_b0 & 320 & Trained (ours) \\
            3 & eca\_nfnet\_l0 & 2304 & Trained (ours) \\
            4 & tf\_efficientnet\_b4 & 448 & External pretrained \\
            5 & regnety\_016 & 888 & External pretrained \\
            6 & regnety\_008 & 768 & External pretrained \\
            \bottomrule
        \end{tabular}
    \end{center}

    \vspace{0.5cm}
    \begin{columns}[T, onlytextwidth]
        \begin{column}{0.46\textwidth}
            \begin{block}{Ensemble Strategy}
                \small
                \textbf{Arithmetic Mean} of softmax:
                $$\hat{p} = \frac{1}{6}\sum_{m=1}^{6} p^{(m)}$$
            \end{block}
        \end{column}
        \begin{column}{0.46\textwidth}
            \begin{alertblock}{Why Diverse Models?}
                \small
                \begin{itemize}
                    \item Different inductive biases
                    \item Error decorrelation
                    \item Robust to distribution shift
                \end{itemize}
            \end{alertblock}
        \end{column}
    \end{columns}
\end{frame}

% ============================================
\section{Results}
% ============================================

% ------------------------------------------
% Performance Results
% ------------------------------------------
\begin{frame}{Performance Results}
    \begin{center}
        \small
        \begin{tabular}{lccc}
            \toprule
            \textbf{Model} & \textbf{Top-1 Acc} & \textbf{Top-5 Acc} & \textbf{vs Random} \\
            \midrule
            Random Baseline & 0.49\% & 2.43\% & 1$\times$ \\
            \midrule
            Base Model (EfficientNet-B3 SED) & 60.07\% & 80.34\% & 122.6$\times$ \\
            \rowcolor{draculacurrent}
            \textbf{6-Model Ensemble} & \textbf{87.99\%} & \textbf{95.99\%} & \textbf{179.6$\times$} \\
            \bottomrule
        \end{tabular}
    \end{center}

    \vspace{0.5cm}
    \begin{columns}[T, onlytextwidth]
        \begin{column}{0.46\textwidth}
            \begin{block}{Key Achievements}
                \small
                \begin{itemize}
                    \item \textbf{+28\%} improvement from ensemble
                    \item \textbf{181$\times$} better than random
                    \item 206 species classified
                    \item Real-time inference capability
                \end{itemize}
            \end{block}
        \end{column}
        \begin{column}{0.46\textwidth}
            \begin{exampleblock}{Technique Contributions}
                \small
                \begin{itemize}
                    \item SED Architecture: Foundation
                    \item Data Augmentation: +5\% acc
                    \item Ensemble: +3\% acc
                    \item TTA + Smoothing: +0.5\% acc
                \end{itemize}
            \end{exampleblock}
        \end{column}
    \end{columns}
\end{frame}

% ------------------------------------------
% Sample Predictions
% ------------------------------------------
\begin{frame}{Sample Audio Predictions}
    \begin{center}
        \small
        \begin{tabular}{lllr}
            \toprule
            \textbf{Audio File} & \textbf{True Species} & \textbf{Top Prediction} & \textbf{Confidence} \\
            \midrule
            XC254589.ogg & amekes & amekes & \textcolor{green!60!black}{\textbf{84.3\%}} \\
            XC493187.ogg & ywcpar & ywcpar & \textcolor{green!60!black}{\textbf{53.0\%}} \\
            XC449442.ogg & savhaw1 & savhaw1 & 22.8\% \\
            XC171396.ogg & smbani & strher & \textcolor{red}{11.7\%} \\
            XC82494.ogg & bobfly1 & crbtan1 & \textcolor{red}{61.6\%} \\
            \bottomrule
        \end{tabular}
    \end{center}

    \vspace{0.2cm}
    \begin{alertblock}{Observations}
        \small
        \begin{itemize}
            \item High confidence for clear recordings (84.3\%)
            \item Confusions occur for acoustically similar species
            \item Rare species remain challenging (class imbalance)
        \end{itemize}
    \end{alertblock}
\end{frame}

% ============================================
\section{Summary}
% ============================================

% ------------------------------------------
% Techniques Summary
% ------------------------------------------
\begin{frame}[t]{All Techniques Implemented}
    \footnotesize
    \begin{multicols}{2}
        \textbf{Architecture:}
        \begin{enumerate}
            \item SED (GeMFreq + AttHead)
            \item 20s chunks, 224 mel bands
            \item Z-score + min-max normalization
            \item Absmax audio normalization
        \end{enumerate}

        \textbf{Training:}
        \begin{enumerate}
            \setcounter{enumi}{4}
            \item Mixup augmentation
            \item SpecAugment
            \item Weighted Random Sampling
            \item Stochastic Depth (DropPath)
            \item Multi-Iterative Noisy Student
            \item CrossEntropy + label smoothing
            \item AdamW + CosineAnnealing
        \end{enumerate}

        \columnbreak

        \textbf{Inference:}
        \begin{enumerate}
            \setcounter{enumi}{11}
            \item Overlap-Average Inference
            \item Delta-Shift TTA
            \item Gaussian Smoothing
            \item Multi-Model Ensemble
        \end{enumerate}

        \vspace{0.2cm}
        \begin{center}
            \textcolor{birdgreen}{\Large \textbf{15 Techniques}}
        \end{center}
    \end{multicols}
\end{frame}

% ------------------------------------------
% Conclusion
% ------------------------------------------
\begin{frame}[t]{Conclusion}
    \centering
    \begin{minipage}{0.6\textwidth}
        \begin{block}{Key Achievements}
            \footnotesize
            \begin{itemize}
                \item Complete SED implementation
                \item State-of-the-art: \textbf{87.99\%}
                \item 6-model diverse ensemble
                \item 15 advanced techniques
                \item Modular, extensible code
            \end{itemize}
        \end{block}
    \end{minipage}

    \vspace{0.5cm}
    \begin{exampleblock}{Real-World Applications}
        \centering \footnotesize Biodiversity Monitoring $\cdot$ Automated Surveys $\cdot$ Climate Change
    \end{exampleblock}
\end{frame}

% ------------------------------------------
% Thank You
% ------------------------------------------
\setbeamercolor{palette primary}{bg=draculapurple, fg=draculabg}
\begin{frame}[standout]
    \Huge Thank You!
    
    \vspace{1cm}
    \normalsize
    Questions?
\end{frame}

\end{document}
